\textbf{Время:} 02:00

\textbf{Тезис:} отношение пар одного события к смешанному контролю оценивает корреляционный избыток.

\textbf{Визуал:} распределения A и B → нормированное отношение; пояснение центральности дать при слове «похожие».

\textbf{Текст:} В A помещаем пары пионов одного события. Его форма зависит и от корреляций, и от спектров и регистрации. Для контроля берём частицы из разных похожих событий и получаем B.

Почему похожих? Центральность — шкала, связанная с геометрией перекрытия ядер. В эксперименте её оценивают, например, по множественности заряженных частиц и калибруют по долям полного адронного сечения. Класс 0–10 процентов — наиболее центральная десятая часть этого сечения, а не десять процентов радиуса ядра. Внутри близких классов смешиваем события с сопоставимыми спектрами. Учитываем также вершину и условия регистрации.

Нормированное отношение A к B оценивает введённое ранее сравнение плотностей выходов частиц и пар. Общая нормировка задаёт C = 1 для независимого фона в контрольной области, где фемтоскопический эффект мал. Все парные эффекты и нефемтоскопические связи смешивание автоматически не устраняет.

На учебном рисунке два счёта сгенерированы независимо по Пуассону, а событий и настоящего смешивания нет. Экспозиция B больше в десять раз; нормировка десять задана заранее. Пример объясняет деление, но не зависимости реального фона.

\textbf{Переход:} Теперь проверим, насколько точно мы знаем полученное отношение.
